% ============================================================
%  Dok 255 | Information als topologisch eingefrorene
%           kinetische Energie
%  Kapitelinhalt DE
% ============================================================

\title{Information als topologisch eingefrorene kinetische Energie\\[0.4em]
\large Landauers Prinzip und Vopsons Äquivalenz als Folgen der \(T^4\)-Geometrie in FFGFT}
\author{Johann Pascher}
\date{Mai 2026}
\maketitle

\begin{abstract}
In der Zeit-Masse-Dualität (FFGFT) gilt \(\tilde{T} \cdot m = 1\),
wobei \(\tilde{T} = 1/m\) die zur Masse duale Zeitvariable auf dem
4D-Identifikationstorus \(T^4\) ist.
Diese Grundrelation erzwingt eine natürliche Zweiteilung der Energie:
statische Energie als topologisch eingeschlossene Kinematik
(Windungszahlen auf \(T^4\)) und freie kinetische Energie
(Bewegung durch \(T^4\)).
Aus dieser Zweiteilung folgt ohne zusätzliche Postulate,
dass Information topologisch eingefrorene kinetische Energie ist.
Landauers Prinzip \cite{Landauer1961} und Vopsons
Masse-Information-Äquivalenz \cite{Vopson2019}
werden als geometrische Konsequenzen der \(T^4\)-Struktur abgeleitet,
nicht als unabhängige Postulate.
Die Quantisierung der Information folgt aus der Diskretheit der Windungszahlen.
Wir arbeiten in natürlichen Einheiten mit \(c = \hbar = k_B = 1\),
sofern nicht anders angegeben.
\end{abstract}

\tableofcontents
\newpage

% ------------------------------------------------------------
\section{Ausgangslage: Zwei Energietypen in \(\tilde{T} \cdot m = 1\)}
% ------------------------------------------------------------

Die Grundrelation der FFGFT \cite{Pascher2026},
\begin{equation}
  \tilde{T} \cdot m = 1,
\end{equation}
verknüpft die zur Masse \(m\) duale Zeitvariable
\(\tilde{T} = 1/m\) mit der Masse als duale Größen auf dem
4D-Identifikationstorus \(T^4\).
Aus dieser Struktur ergeben sich zwei qualitativ verschiedene
Energieformen:

\paragraph{Statische Energie.}
An Ruhemasse gebundene Energie \(E_0 = mc^2\).
Sie entspricht topologisch stabilen Windungszahlen auf \(T^4\) --
geometrischen Invarianten, die ohne Wechselwirkung erhalten bleiben.
Diese Energie ist nicht statistisch; sie ist ein diskreter topologischer
Zustand. Sie trägt nicht zur thermodynamischen Entropie bei,
solange die Topologie intakt bleibt.

\paragraph{Dynamische kinetische Energie.}
Freie Bewegungsenergie -- bei masselosen Teilchen (Photonen) die
einzige Energieform \(E = hf\), bei massebehafteten Teilchen
der Anteil über die Ruhemasse hinaus.
Diese Energie ist statistisch verteilbar und steht mit
thermodynamischen Systemen im Gleichgewicht.

Das \textbf{Photon} bildet den Grenzfall:
ohne Ruhemasse, ohne Windungszahl für Masse,
rein kinetisch, rein statistisch.
Es ist der natürliche Träger von Shannon-Information
im Sinne von Wheeler \cite{Wheeler1990}.

% ------------------------------------------------------------
\section{Masse als topologisch eingefrorene Kinematik}
% ------------------------------------------------------------

\subsection{Die Grundthese}

Vor jeder Massenentstehung existiert ausschließlich kinetische Energie.
Masse entsteht, wenn kinetische Energie durch die \(T^4\)-Topologie
in eine stabile Windungsstruktur eingeschlossen wird.

\begin{center}
\emph{Masse ist topologisch eingefrorene kinetische Energie.}
\end{center}

\paragraph{Definition der Windungszahl.}
Sei \(\gamma\) ein geschlossener Pfad auf \(T^4\).
Die Windungszahl \(w \in \mathbb{Z}\) ist die ganzzahlige
Homotopieklasse \([\gamma] \in \pi_1(T^4) \cong \mathbb{Z}^4\).
Masse entspricht einer stabilen Kombination solcher Windungszahlen.
In FFGFT tragen alle Leptonen spezifische Windungsexponenten \(p_i\),
die aus der Geometrie von \(T^4\) folgen \cite{Pascher2026}.

\subsection{Visualisierung: Freie und eingefrorene Kinematik}

\begin{figure}[htbp]
\centering
\begin{tikzpicture}[
  node distance=1.2cm,
  every node/.style={font=\small},
  box/.style={draw, rounded corners=4pt, minimum width=2.8cm,
              minimum height=1.0cm, align=center, fill=blue!8},
  massbox/.style={draw, rounded corners=4pt, minimum width=2.8cm,
                  minimum height=1.0cm, align=center, fill=orange!15},
  arr/.style={-{Stealth[length=6pt]}, thick},
  redarr/.style={-{Stealth[length=6pt]}, thick, red!70!black},
]

% Nodes
\node[box] (kin) {Freie kinetische\\Energie\\(statistisch)};
\node[massbox, right=2.2cm of kin] (top) {Topologischer\\Einschluss\\(Windungszahl \(w\))};
\node[massbox, right=2.2cm of top] (mass) {Ruhemasse /\\Information\\(statisch)};

% Forward arrows
\draw[arr] (kin) -- node[above,font=\footnotesize]{Einfrieren} (top);
\draw[arr] (top) -- node[above,font=\footnotesize]{Speichern} (mass);

% Reverse arrow (erasure)
\draw[redarr, bend right=35] (mass) to
  node[below,font=\footnotesize,red!70!black]{Löschen: \(E_L = k_BT\ln 2\)}
  (kin);

% Labels below
\node[below=0.5cm of kin, font=\footnotesize, text=blue!70!black]
  {Photon: immer hier};
\node[below=0.5cm of mass, font=\footnotesize, text=orange!80!black]
  {Elektron: hier + links};

\end{tikzpicture}
\caption{Zweiteilung der Energie in FFGFT.
  Freie kinetische Energie (blau) kann durch topologischen Einschluss
  in statische Masse/Information (orange) überführt werden.
  Das Löschen (roter Pfeil) setzt die eingefrorene Energie als
  Wärme frei -- Landauers Prinzip.
  Das Photon verbleibt dauerhaft im kinetischen Sektor.}
\label{fig:energie}
\end{figure}

\subsection{Analogie: Stehende Welle}

Eine stehende Welle auf einer Saite ist kinetisch --
die Saite schwingt. Dennoch erscheint das Muster statisch.
Analog ist die Ruhemasse eines Elektrons die stehende
Welle einer topologisch eingeschlossenen Schwingung auf \(T^4\).
Die \enquote{Stille} der Ruhemasse ist die Stille des Musters,
nicht der Substanz.

\subsection{Der Higgs-Mechanismus in diesem Bild}

Im Standardmodell verleiht der Higgs-Mechanismus den Teilchen Masse
durch Symmetriebrechung. In der \(T^4\)-Sprache entspricht
dies der Selektion stabiler topologischer Sektoren: kinetische Energie
wird in Windungsstrukturen eingeschlossen.
Die FFGFT-Gleichung \(\xi = \lambda_h^2 v^2 / (16\pi^3 m_h^2)\)
ist der algebraische Ausdruck dieses geometrischen
Selektionsprozesses -- nicht seine Ersetzung.

% ------------------------------------------------------------
\section{Information als topologischer Einschluss}
% ------------------------------------------------------------

Wenn Masse topologisch eingefrorene kinetische Energie ist,
dann gilt für Information:

\begin{center}
\renewcommand{\arraystretch}{1.4}
\begin{tabular}{lll}
\toprule
\textbf{Operation} & \textbf{Geometrisch} & \textbf{Energetisch} \\
\midrule
Erzeugen  & Windungszahl \(w \to w\pm1\) & Kinetisch einfrieren \\
Speichern & \(w = \mathrm{const}\)       & Topologie erhalten \\
Löschen   & \(w \to 0\)                  & \(E_L = k_BT\ln2\) freigeben \\
\bottomrule
\end{tabular}
\end{center}

\paragraph{Definition des Informationsgehalts.}
Der Informationsgehalt \(I(\gamma)\) eines topologischen Zustands \(\gamma\)
ist die minimale Anzahl von Bits, die nötig ist,
um seine Windungszahlklasse \([\gamma] \in \mathbb{Z}^4\) zu spezifizieren.
\(I\) ist ganzzahlig.
Shannon-Entropie entsteht als statistischer Erwartungswert über
Ensembles solcher Zustände -- analog zur Temperatur als Mittelwert
über diskrete kinetische Energien.

% ------------------------------------------------------------
\section{Landauers Prinzip aus \(T^4\)-Geometrie}
% ------------------------------------------------------------

\subsection{Das klassische Prinzip}

Landauer (1961) zeigt, dass das Löschen eines Bits
irreversibel Wärme an die Umgebung abgibt \cite{Landauer1961}:
\begin{equation}
  E_L = k_B T \ln 2 \quad \text{pro Bit.}
  \label{eq:landauer}
\end{equation}
Dies wurde experimentell von Bérut et al.\ (2012) mit einem
kolloidalen Teilchen in einem Doppelmuldenpotential bestätigt
\cite{Berut2012}.

\subsection{Geometrische Herleitung}

Sei eine minimale topologische Unterscheidung auf \(T^4\) gegeben
durch zwei Zustände mit Windungszahldifferenz \(\Delta w = 1\).
Die zugehörige eingefrorene kinetische Energie ist
\(E_0 = \Delta w \cdot \epsilon_0\) mit einem fundamentalen Quant \(\epsilon_0\).

Im thermischen Gleichgewicht bei Temperatur \(T\) beschreibt der
Boltzmann-Faktor die Wahrscheinlichkeit, diese Unterscheidung aufzulösen.
Die dabei freigesetzte Energie ist nicht kleiner als \(k_B T \ln 2\),
da dies die minimale Arbeit ist, um bei Temperatur \(T\) eine binäre
Unterscheidung irreversibel zu löschen (aus dem zweiten Hauptsatz).

Die \(T^4\)-Topologie erzwingt nun, dass
\(\epsilon_0 = k_B T \ln 2\):
die eingefrorene kinetische Energie einer minimalen Windungsdifferenz
ist derselbe physikalische Betrag wie die minimale Löscharbeit aus
dem zweiten Hauptsatz -- zwei Ausdrücke für dieselbe Größe,
keine neue Annahme. Damit ist Landauers Prinzip eine
geometrische Konsequenz -- kein unabhängiges thermodynamisches Postulat.

% ------------------------------------------------------------
\section{Vopsons Masse-Information-Äquivalenz als Konsequenz}
% ------------------------------------------------------------

Vopson (2019) leitet empirisch aus Landauers Prinzip die
Masse-Information-Äquivalenz ab \cite{Vopson2019}:
\begin{equation}
  M_I = \frac{k_B T \ln 2}{c^2} \quad \text{pro Bit.}
  \label{eq:vopson}
\end{equation}
In der \(T^4\)-Sprache ist \(M_I\) die Masse,
die einer minimalen topologischen Windungszahldifferenz
\(\Delta w = 1\) bei Temperatur \(T\) entspricht.

\textbf{Vopson entdeckte empirisch, was FFGFT geometrisch erklärt.}
Die Ableitungshierarchie ist:
\begin{equation}
  T^4\text{-Topologie} \;\longrightarrow\;
  \Delta w = 1 \;\longrightarrow\;
  E_L = k_B T \ln 2 \;\longrightarrow\;
  M_I = E_L / c^2.
\end{equation}

% ------------------------------------------------------------
\section{Quantisierung der Information}
% ------------------------------------------------------------

Windungszahlen auf \(T^4\) sind diskrete ganze Zahlen
(\(\pi_1(T^4) \cong \mathbb{Z}^4\)).
Daraus folgt:

\begin{itemize}
  \item Information ist natürlich quantisiert --
        nicht als Konvention, sondern als topologische Notwendigkeit.
  \item Die minimale Informationseinheit ist \(\Delta w = 1\)
        (eine Windungsänderung in einer der vier Torus-Richtungen).
  \item Die räumliche Skala dieser minimalen Einheit ist die
        fundamentale Länge der FFGFT:
        \begin{equation}
          L_0 = \xi \cdot l_P \approx 2{,}15 \times 10^{-39}\,\text{m},
        \end{equation}
        wobei \(\xi = 4/3 \times 10^{-4}\) der universelle geometrische
        Parameter und \(l_P\) die Planck-Länge ist.
        \(L_0\) ist damit die kleinste physikalisch bedeutsame
        Körnigkeit der Information in FFGFT -- unterhalb dieser
        Skala ist keine topologische Unterscheidung möglich.
  \item Kontinuierliche Shannon-Entropie entsteht als statistisches
        Mittel über diskrete Topologiezustände.
\end{itemize}

% ------------------------------------------------------------
\section{Das Photon als Grenzfall}
% ------------------------------------------------------------

Das Photon ist der einzige masselose Zustand im FFGFT-Spektrum,
der keine topologische Windungsstruktur für Masse besitzt.
Daher ist es rein kinetisch und verbleibt dauerhaft im statistischen
Sektor (vgl.\ Abb.~\ref{fig:energie}).

\begin{itemize}
  \item Das Photon ist der natürliche Träger von Shannon-Information
        (\enquote{It from Bit} \cite{Wheeler1990}).
  \item Photonische Information ist nicht einfrierbar
        in Ruhemasse ohne Wechselwirkung mit massebehafteter Materie.
  \item Die Lichtgeschwindigkeit \(c\) begrenzt die
        Übertragungsrate freier kinetischer Information --
        als geometrische Konsequenz von \(T^4\), nicht als Postulat.
\end{itemize}

% ------------------------------------------------------------
\section{Verbindung zur Unterscheidbarkeitsstruktur nach Haskins}
% ------------------------------------------------------------

Haskins (2026, eingereicht bei Foundations of Physics \cite{Haskins2026})
beschreibt eine Emergenz-Hierarchie:
\begin{align*}
C_a\Psi = 0 \;&\to\; \text{Multiplizität} \;\to\;
\text{Unterscheidbarkeitskollaps} \\
&\to\; \text{Äquivalenzstruktur} \;\to\;
\text{Grobkörnungsprojektion} \\
&\to\; \text{Entartungsstruktur} \;\to\;
\text{strukturelle Entropieordnung.}
\end{align*}

In der \(T^4\)-Sprache kann der Unterscheidbarkeitskollaps als
der Moment interpretiert werden, in dem zwei kinetische Zustände durch
topologischen Einschluss in dieselbe Windungsklasse fallen --
also nicht mehr unterscheidbar sind.
Die Constraint-Bedingung \(C_a\Psi = 0\) selektiert dann die
topologisch stabilen Sektoren auf \(T^4\).

\textbf{Diese Verbindung ist eine interpretative Brücke, keine Ableitung.}
Eine formale Entsprechung zwischen Haskins' Constraints und
den \(T^4\)-Windungszahlen ist Gegenstand weiterer Arbeit.

% ------------------------------------------------------------
\section{Zusammenfassung}
% ------------------------------------------------------------

\begin{enumerate}
  \item \(\tilde{T} \cdot m = 1\) erzwingt zwei Energietypen:
        topologisch eingeschlossen (statisch) und frei (kinetisch).
  \item Masse ist topologisch eingefrorene kinetische Energie.
  \item Information ist eingefrorene kinetische Energie
        in topologischer Form (Windungszahl \(\Delta w = 1\) = 1 Bit).
  \item Landauers Prinzip folgt geometrisch aus dem Boltzmann-Faktor
        für Windungsauflösung; \(E_L = k_BT\ln 2\) ist kein Postulat.
  \item Vopsons Äquivalenz \(M_I = k_BT\ln2/c^2\) ist eine direkte
        Konsequenz -- Vopson entdeckte empirisch, was FFGFT erklärt.
  \item Quantisierung der Information folgt aus \(\pi_1(T^4)\cong\mathbb{Z}^4\).
  \item Das Photon ist der Grenzfall: rein kinetisch, rein statistisch,
        natürlicher Informationsträger.
\end{enumerate}

\bigskip
\noindent
\textbf{Methodologische Einordnung:}
Diese Ableitung ist eine Brücke (algebraisch aus \(\tilde{T}\cdot m=1\)
und \(T^4\)-Topologie).
Was sie leistet: Landauer und Vopson sind keine unabhängigen Postulate.
Was sie nicht leistet: eine experimentelle Unterscheidung gegenüber
anderen Theorien, die dieselben Gleichungen reproduzieren.
Dazu wären spezifische Vorhersagen nötig (z.\,B.\ Abweichungen bei
sehr tiefen Temperaturen oder in starken Gravitationsfeldern).

\bigskip
\begin{flushright}
\textit{Johann Pascher, Oberösterreich, Mai 2026}
\end{flushright}

% ------------------------------------------------------------
\begin{thebibliography}{99}

\bibitem{Landauer1961}
R.~Landauer,
\textit{Irreversibility and heat generation in the computing process},
IBM Journal of Research and Development \textbf{5}(3), 183--191 (1961).

\bibitem{Berut2012}
A.~Bérut, A.~Arakelyan, A.~Petrosyan, S.~Ciliberto,
R.~Dillenschneider, E.~Lutz,
\textit{Experimental verification of Landauer's principle linking
information and thermodynamics},
Nature \textbf{483}, 187--189 (2012).

\bibitem{Vopson2019}
M.\,M.~Vopson,
\textit{The mass-energy-information equivalence principle},
AIP Advances \textbf{9}(9), 095206 (2019).

\bibitem{Wheeler1990}
J.\,A.~Wheeler,
\textit{Information, physics, quantum: The search for links},
in W.~Zurek (ed.), \textit{Complexity, Entropy, and the Physics
of Information}, Addison-Wesley, 3--28 (1990).

\bibitem{Haskins2026}
D.~Haskins,
\textit{Admissibility-constrained distinguishability and the emergence
of coarse-grained structure},
Foundations of Physics (eingereicht, 2026).

\bibitem{Pascher2026}
J.~Pascher,
\textit{Fundamentale Fraktale Geometrische Feldtheorie (FFGFT) --
T0-Zeit-Masse-Dualität},
Zenodo, DOI: 10.5281/zenodo.20117635 (2026).

\end{thebibliography}
